% Options for packages loaded elsewhere
\PassOptionsToPackage{unicode}{hyperref}
\PassOptionsToPackage{hyphens}{url}
\documentclass[
]{article}
\usepackage{xcolor}
\usepackage{amsmath,amssymb}
\setcounter{secnumdepth}{-\maxdimen} % remove section numbering
\usepackage{iftex}
\ifPDFTeX
  \usepackage[T1]{fontenc}
  \usepackage[utf8]{inputenc}
  \usepackage{textcomp} % provide euro and other symbols
\else % if luatex or xetex
  \usepackage{unicode-math} % this also loads fontspec
  \defaultfontfeatures{Scale=MatchLowercase}
  \defaultfontfeatures[\rmfamily]{Ligatures=TeX,Scale=1}
\fi
\usepackage{lmodern}
\ifPDFTeX\else
  % xetex/luatex font selection
\fi
% Use upquote if available, for straight quotes in verbatim environments
\IfFileExists{upquote.sty}{\usepackage{upquote}}{}
\IfFileExists{microtype.sty}{% use microtype if available
  \usepackage[]{microtype}
  \UseMicrotypeSet[protrusion]{basicmath} % disable protrusion for tt fonts
}{}
\makeatletter
\@ifundefined{KOMAClassName}{% if non-KOMA class
  \IfFileExists{parskip.sty}{%
    \usepackage{parskip}
  }{% else
    \setlength{\parindent}{0pt}
    \setlength{\parskip}{6pt plus 2pt minus 1pt}}
}{% if KOMA class
  \KOMAoptions{parskip=half}}
\makeatother
\usepackage{longtable,booktabs,array}
\usepackage{calc} % for calculating minipage widths
% Correct order of tables after \paragraph or \subparagraph
\usepackage{etoolbox}
\makeatletter
\patchcmd\longtable{\par}{\if@noskipsec\mbox{}\fi\par}{}{}
\makeatother
% Allow footnotes in longtable head/foot
\IfFileExists{footnotehyper.sty}{\usepackage{footnotehyper}}{\usepackage{footnote}}
\makesavenoteenv{longtable}
\setlength{\emergencystretch}{3em} % prevent overfull lines
\providecommand{\tightlist}{%
  \setlength{\itemsep}{0pt}\setlength{\parskip}{0pt}}
\usepackage{bookmark}
\IfFileExists{xurl.sty}{\usepackage{xurl}}{} % add URL line breaks if available
\urlstyle{same}
\hypersetup{
  hidelinks,
  pdfcreator={LaTeX via pandoc}}

\author{}
\date{}

\begin{document}

\section{VariableNode Resolution
Specification}\label{variablenode-resolution-specification}

\subsection{Variable Kinds}\label{variable-kinds}

There are exactly five kinds of variable a \texttt{VariableNode} can
represent. The resolution logic tries them in priority order; first
match wins.

\subsubsection{1. Context Variable}\label{context-variable}

A variable whose name matches a key in
\texttt{expression.context.variables}. These are externally provided
named constants or parameters.

\begin{itemize}
\tightlist
\item
  \textbf{Resolved by:} \texttt{resolveContextualVariable()}
\item
  \textbf{\texttt{reference.type} set to:}
  \texttt{instanceVariable.getClass()}
\item
  \textbf{Codegen:} \texttt{loadThisAndFieldOntoStack} (field access on
  generated class)
\item
  \textbf{Examples:}

  \begin{itemize}
  \tightlist
  \item
    \texttt{a} in \texttt{RealFunction.express("x+a",\ context)} where
    \texttt{context.registerVariable("a",\ someReal)}
  \item
    \texttt{v}, \texttt{z} in Bessel expressions where they are
    context-provided parameters
  \item
    \texttt{α}, \texttt{β} in
    \texttt{P:n➔when(n=0,1,n=1,(C(1)*x-β+α)/2,...)} (Jacobi polynomials)
  \end{itemize}
\end{itemize}

\subsubsection{2. Independent Variable
(Non-Nullary)}\label{independent-variable-non-nullary}

The input parameter of the function. Either explicitly named via arrow
syntax (\texttt{n➔...}, \texttt{x➔...}, \texttt{t➔...}) or the first
free variable encountered when no arrow is present. Only exists when
\texttt{expression.isNullaryFunction()} is false.

\begin{itemize}
\tightlist
\item
  \textbf{Resolved by:} \texttt{isIndependent()} returns true,
  \texttt{!expression.isNullaryFunction()}
\item
  \textbf{\texttt{reference.type} set to:}
  \texttt{expression.domainType}
\item
  \textbf{Codegen:} \texttt{loadInputParameter} (first argument of
  \texttt{evaluate()})
\item
  \textbf{Examples:}

  \begin{itemize}
  \tightlist
  \item
    \texttt{x} in \texttt{RealFunction.express("x\^{}2+1")} --- domain
    is \texttt{Real}
  \item
    \texttt{n} in
    \texttt{RationalFunctionSequence.express("n➔(½-z/2)ⁿ")} --- domain
    is \texttt{Integer}
  \item
    \texttt{t} in \texttt{RealFunction.express("sin(t)")}
  \end{itemize}
\end{itemize}

\subsubsection{3. Independent Variable
(Nullary)}\label{independent-variable-nullary}

When the expression is nullary (\texttt{domainType\ ==\ Object.class}),
there is no input parameter, but a variable still appears. This variable
is the placeholder variable of the codomain type. Since there is no
input parameter, the independent variable and the placeholder variable
are the same thing.

\begin{itemize}
\tightlist
\item
  \textbf{Resolved by:} \texttt{isIndependent()} returns true,
  \texttt{expression.isNullaryFunction()}
\item
  \textbf{\texttt{reference.type} set to:}
  \texttt{expression.coDomainType}
\item
  \textbf{Codegen:} \texttt{generateFunctionalVariableIdentity}
  (allocate + \texttt{.identity()})
\item
  \textbf{Examples:}

  \begin{itemize}
  \tightlist
  \item
    \texttt{x} in \texttt{RationalFunction.express("x/2")} --- codomain
    is \texttt{RationalFunction}
  \item
    \texttt{x} in
    \texttt{RealPolynomialNullaryFunction.express("x\^{}2+1")} ---
    codomain is \texttt{RealPolynomial}
  \item
    \texttt{x} in \texttt{RationalFunction.express("½-x/2")} → produces
    \texttt{(-x+1)/2}
  \end{itemize}
\end{itemize}

\subsubsection{4. Placeholder Variable
(NEW)}\label{placeholder-variable-new}

When the expression is non-nullary AND the codomain implements
\texttt{Function} (i.e.~\texttt{expression.isInterfaceFunctional()} is
true), a second free variable may appear that is neither the independent
variable nor a context variable nor an upstream variable. This is the
\textbf{placeholder variable} of the codomain type --- the variable of
the output polynomial/rational function/etc.

Only \textbf{one} placeholder variable is allowed per expression. A
second unresolved free variable throws \texttt{CompilerException}.

\begin{itemize}
\tightlist
\item
  \textbf{Resolved by:} new resolution step: not context, not
  independent, not upstream, \texttt{expression.isInterfaceFunctional()}
  is true, \texttt{expression.placeholderVariable\ ==\ null}
\item
  \textbf{\texttt{reference.type} set to:}
  \texttt{expression.coDomainType}
\item
  \textbf{Codegen:} \texttt{generateFunctionalVariableIdentity}
  (allocate + \texttt{.identity()})
\item
  \textbf{New field on Expression:}
  \texttt{public\ VariableNode\textless{}D,C,F\textgreater{}\ placeholderVariable}
\item
  \textbf{Examples:}

  \begin{itemize}
  \tightlist
  \item
    \texttt{z} in
    \texttt{RationalFunctionSequence.express("n➔(½-z/2)ⁿ")} ---
    \texttt{n} is independent (Integer), \texttt{z} is placeholder
    (RationalFunction)
  \item
    \texttt{x} in \texttt{P:n➔when(n=0,1,n=1,(C(1)*x-β+α)/2,else,...)}
    --- \texttt{n} is independent, \texttt{α}/\texttt{β} are context,
    \texttt{x} is placeholder
  \item
    \texttt{x} in \texttt{f:k➔√((2*k+½)/π)*...*j(2*k,x)} --- \texttt{k}
    is independent, \texttt{x} is placeholder
  \end{itemize}
\end{itemize}

\subsubsection{5. Upstream Variable}\label{upstream-variable}

A variable that matches the independent variable of a containing
(upstream) expression. Occurs inside sub-expressions created by Σ, Π, or
lambda bodies.

\begin{itemize}
\tightlist
\item
  \textbf{Resolved by:} \texttt{resolveUpstreamVariables()} walks
  \texttt{expression.acceptUntil()}
\item
  \textbf{\texttt{reference.type} set to:} upstream expression's
  \texttt{domainType}
\item
  \textbf{\texttt{upstreamInput} set to:} \texttt{true}
\item
  \textbf{Codegen:} linked field from containing class
\item
  \textbf{Examples:}

  \begin{itemize}
  \tightlist
  \item
    \texttt{x} inside a product body \texttt{Πk➔f(k,x)\{k=1…N\}} where
    \texttt{x} is the outer expression's input
  \item
    \texttt{t} in nested lambda \texttt{n➔t➔sin(tⁿ)} --- \texttt{t} is
    the inner expression's independent variable, \texttt{n} is upstream
  \end{itemize}
\end{itemize}

\subsubsection{6. Undefined}\label{undefined}

None of the above. Throws \texttt{CompilerException}.

\begin{center}\rule{0.5\linewidth}{0.5pt}\end{center}

\subsection{Resolution Priority Table}\label{resolution-priority-table}

\begin{longtable}[]{@{}
  >{\raggedright\arraybackslash}p{(\linewidth - 8\tabcolsep) * \real{0.1818}}
  >{\raggedright\arraybackslash}p{(\linewidth - 8\tabcolsep) * \real{0.2000}}
  >{\raggedright\arraybackslash}p{(\linewidth - 8\tabcolsep) * \real{0.1091}}
  >{\raggedright\arraybackslash}p{(\linewidth - 8\tabcolsep) * \real{0.3455}}
  >{\raggedright\arraybackslash}p{(\linewidth - 8\tabcolsep) * \real{0.1636}}@{}}
\toprule\noalign{}
\begin{minipage}[b]{\linewidth}\raggedright
Priority
\end{minipage} & \begin{minipage}[b]{\linewidth}\raggedright
Condition
\end{minipage} & \begin{minipage}[b]{\linewidth}\raggedright
Kind
\end{minipage} & \begin{minipage}[b]{\linewidth}\raggedright
\texttt{reference.type}
\end{minipage} & \begin{minipage}[b]{\linewidth}\raggedright
Codegen
\end{minipage} \\
\midrule\noalign{}
\endhead
\bottomrule\noalign{}
\endlastfoot
1 & Name found in \texttt{context.variables} & Context &
\texttt{instanceVariable.getClass()} & field access \\
2 & \texttt{isIndependent()} AND nullary & Independent (nullary) &
\texttt{coDomainType} & \texttt{.identity()} \\
3 & \texttt{isIndependent()} AND NOT nullary & Independent (input) &
\texttt{domainType} & \texttt{loadInputParameter} \\
4 & \texttt{isInterfaceFunctional()} AND
\texttt{placeholderVariable\ ==\ null} & Placeholder &
\texttt{coDomainType} & \texttt{.identity()} \\
5 & Found in upstream expression & Upstream & upstream's
\texttt{domainType} & linked field \\
6 & None & UNDEFINED & --- & throw \texttt{CompilerException} \\
\end{longtable}

\textbf{Row 4 is the only new addition.} It requires a new
\texttt{placeholderVariable} field on \texttt{Expression}. The field is
assigned during resolution. If a second variable hits row 4 and
\texttt{placeholderVariable} is already assigned, it falls through to
row 6 and throws --- only one placeholder is allowed.

\begin{center}\rule{0.5\linewidth}{0.5pt}\end{center}

\subsection{Relationship Between Independent and
Placeholder}\label{relationship-between-independent-and-placeholder}

\begin{longtable}[]{@{}
  >{\raggedright\arraybackslash}p{(\linewidth - 10\tabcolsep) * \real{0.1345}}
  >{\raggedright\arraybackslash}p{(\linewidth - 10\tabcolsep) * \real{0.1849}}
  >{\raggedright\arraybackslash}p{(\linewidth - 10\tabcolsep) * \real{0.2185}}
  >{\raggedright\arraybackslash}p{(\linewidth - 10\tabcolsep) * \real{0.1765}}
  >{\raggedright\arraybackslash}p{(\linewidth - 10\tabcolsep) * \real{0.1765}}
  >{\raggedright\arraybackslash}p{(\linewidth - 10\tabcolsep) * \real{0.1092}}@{}}
\toprule\noalign{}
\begin{minipage}[b]{\linewidth}\raggedright
Expression Type
\end{minipage} & \begin{minipage}[b]{\linewidth}\raggedright
\texttt{isNullaryFunction()}
\end{minipage} & \begin{minipage}[b]{\linewidth}\raggedright
\texttt{isInterfaceFunctional()}
\end{minipage} & \begin{minipage}[b]{\linewidth}\raggedright
Independent Variable
\end{minipage} & \begin{minipage}[b]{\linewidth}\raggedright
Placeholder Variable
\end{minipage} & \begin{minipage}[b]{\linewidth}\raggedright
Relationship
\end{minipage} \\
\midrule\noalign{}
\endhead
\bottomrule\noalign{}
\endlastfoot
\texttt{RealFunction} (Real→Real) & false & false & \texttt{x}
(domainType=Real) & N/A & No placeholder possible; codomain is scalar \\
\texttt{RealNullaryFunction} (Object→Real) & true & false & \texttt{x}
if present (coDomainType=Real) & N/A & No placeholder possible; codomain
is scalar \\
\texttt{RationalNullaryFunction} (Object→RationalFunction) & true & true
& \texttt{x} (coDomainType=RationalFunction) & same as independent &
They collapse to one variable because there's no input \\
\texttt{RationalFunctionSequence} (Integer→RationalFunction) & false &
true & \texttt{n} (domainType=Integer) & \texttt{z}
(coDomainType=RationalFunction) & Distinct variables; \texttt{n} is
input, \texttt{z} is placeholder \\
\texttt{RealPolynomialFunction} (RealPolynomial→RealPolynomial) & false
& true & \texttt{x} (domainType=RealPolynomial) & could exist if a
second free var appears & Independent is the input polynomial;
placeholder would be the output polynomial's variable \\
\texttt{RealBivariateFunction} (Real→RealFunction) & false & true
(codomain is RealFunction) & first variable (domainType=Real) & second
variable (coDomainType=RealFunction) & Classic bivariate split \\
\end{longtable}

\begin{center}\rule{0.5\linewidth}{0.5pt}\end{center}

\subsection{Why This Is Complete}\label{why-this-is-complete}

Every variable encountered during parsing must be exactly one of:

\begin{enumerate}
\def\labelenumi{\arabic{enumi}.}
\tightlist
\item
  \textbf{Already known} --- its name is in \texttt{context.variables} →
  Context Variable
\item
  \textbf{The input} --- it matches or becomes the independent variable
  → Independent Variable
\item
  \textbf{The output type's variable} --- the codomain is functional and
  no placeholder is yet assigned → Placeholder Variable
\item
  \textbf{From an enclosing scope} --- it matches an upstream
  expression's independent variable → Upstream Variable
\item
  \textbf{None of the above} → Undefined, throw error
\end{enumerate}

There is no sixth option. There is no stack of placeholder variables.
There is no concept of ``formal'', ``declared'', or ``indeterminate''
variable. A placeholder variable is simply the single variable that the
output functional type uses as its own variable, declared by its
appearance in the expression, limited to exactly one per expression.

\end{document}
